%% 17장 배열과 표의 메모리 모델
\chapter{배열과 표의 메모리 모델}
\label{ch:memory-model}

\abstract{수치 계산의 데이터는 거의 언제나 다차원 배열 아니면 열로 된 표이다. 이 장에서는 두 구조가 메모리에 어떻게 놓이는지를 다룬다. 다차원 배열은 연속된 메모리 블록 하나와 그것을 읽는 규칙인 모양과 보폭으로 이루어지는데, 이 단순한 표현 덕분에 전치와 잘라 내기와 브로드캐스팅이 메모리를 전혀 옮기지 않는 연산이 된다. 그 위에서 원소별 연산을 하는 반복자가 축을 다시 배열하고 합쳐 긴 연속 구간을 만드는 방법, 서로 다른 자료형을 섞어 계산할 때의 승격 규칙, 입력과 출력이 메모리를 공유할 때의 처리를 본다. 이어서 데이터프레임의 열이 값, 결측을 나타내는 비트맵, 문자열의 위치표로 이루어지는 열 지향 형식과, 복사를 미루었다가 쓰기가 일어날 때만 하는 Copy-on-Write를 다룬다.}

\section{보폭으로 나타낸 배열}
\label{sec:strides}

\subsection{모양과 보폭}

$d$차원 배열은 \term{모양}{shape} $(n_0,\dots,n_{d-1})$과 \term{보폭}{stride} $(s_0,\dots,s_{d-1})$, 그리고 시작 주소 $p$로 나타낸다. 첨자 $(i_0,\dots,i_{d-1})$의 원소는 주소
\[
p+\sum_{k=0}^{d-1}i_ks_k
\]
에 있다. 보폭은 바이트 단위이고, 첨자에서 주소로 가는 이 사상은 아핀 사상이다. 원소 하나가 $b$바이트일 때 행 우선, 곧 C 순서로 빈틈없이 저장한 배열의 보폭은 $s_{d-1}=b$, $s_k=n_{k+1}s_{k+1}$이고, 열 우선, 곧 Fortran 순서의 보폭은 $s_0=b$, $s_k=n_{k-1}s_{k-1}$이다. 예컨대 모양 $(3,4,5)$인 binary64 배열의 보폭은 C 순서에서 $(160,40,8)$, F 순서에서 $(8,24,96)$이다.

이 표현의 힘은 보폭을 바꾸는 것만으로 많은 연산을 할 수 있다는 데 있다. 두 축을 바꾸는 전치는 모양과 보폭에서 두 성분의 자리를 바꾸는 것이고, $a[1{:}7{:}2]$ 같은 잘라 내기는 시작 주소를 $s_0$만큼 옮기고 보폭을 두 배로 하는 것이며, 축을 뒤집는 것은 보폭의 부호를 바꾸는 것이다. 어느 경우에도 원소는 한 바이트도 옮겨지지 않으므로, $1$GB 배열의 전치와 $2\times2$ 배열의 전치는 비용이 같다. 이렇게 같은 메모리를 다른 규칙으로 읽는 배열을 \term{뷰}{view}라 하며, numpy가 이 표현을 Python 과학 계산의 공통 기반으로 만들었다\cite{vanderwalt2011numpy,harris2020numpy}.

\subsection{연속성}

어떤 연산은 배열이 메모리에 빈틈없이 한 줄로 놓여 있을 때만, 또는 그때 훨씬 빠르게 할 수 있다. 보폭이 C 순서의 공식과 같으면 \term{C 연속}{C-contiguous}, F 순서의 공식과 같으면 \term{F 연속}{F-contiguous}이라 한다. 여기서 크기가 $1$인 축의 보폭은 어떤 값이든 주소에 영향을 주지 않으므로 비교에서 빼고, 원소가 없는 배열은 양쪽 모두 연속으로 본다. 이것이 numpy가 쓰는 너그러운 규칙이다.

모양을 바꾸는 \code{reshape}이 뷰가 될 수 있는지도 보폭으로 판정한다. 원소의 순서를 C 순서로 읽었을 때 바뀌지 않는다면, 새 모양의 보폭을 옛 보폭에서 만들 수 있을 때만 뷰가 된다. 연속인 배열은 언제나 그렇지만, 전치된 배열을 평평하게 만들 때처럼 C 순서로 읽은 원소들이 메모리에서 일정한 간격으로 놓여 있지 않으면 보폭으로 나타낼 수 없으므로 복사해야 한다.

\subsection{브로드캐스팅}

모양이 다른 두 배열의 원소별 연산을 정의하는 규칙이 \term{브로드캐스팅}{broadcasting}이다. 두 모양을 오른쪽 끝에 맞추어 늘어놓고 짧은 쪽의 왼쪽에 크기 $1$인 축을 채운 뒤, 축마다 두 크기가 같거나 하나가 $1$이어야 한다. 결과의 크기는 둘 가운데 큰 쪽이다. 크기 $1$인 축을 늘리는 것은 그 축의 보폭을 $0$으로 두는 것과 같다. 보폭이 $0$이면 그 축의 첨자가 무엇이든 같은 원소를 가리키기 때문이다. 그래서 모양 $(3,1)$인 binary64 배열과 모양 $(4,)$인 배열을 더하면 결과의 모양은 $(3,4)$이고, 두 입력은 각각 보폭 $(8,0)$과 $(0,8)$로 읽힌다. 브로드캐스팅도 메모리를 옮기지 않는다.

보폭이 $0$인 뷰에는 여러 첨자가 같은 메모리를 가리키므로 그런 뷰에 쓰기를 허용하면 위험하다. 한 원소에 쓴 값이 다른 원소에도 나타나기 때문이다. 그래서 브로드캐스트한 뷰는 읽기 전용으로 만든다.

\begin{instardust}[\code{docs/design/array.md} §3.3]
stardust의 \code{NDArray}는 메모리를 소유하는 참조 계수 버퍼, 그 안의 시작 주소, 자료형, 그리고 모양과 보폭으로 이루어진다. 모양과 보폭은 객체 안에 크기 $32$의 고정 배열로 두어 배열을 만들 때마다 힙 할당을 하지 않으므로, 차원의 수는 $32$를 넘을 수 없다. 객체를 복사할 때는 고정 배열 전체가 아니라 앞의 \code{ndim}칸만 옮기는데, 설계 문서는 $512$바이트를 통째로 복사하는 것만으로도 작은 연산의 비용이 눈에 띄게 늘었다고 적는다. 연속성 플래그는 모양과 보폭이 바뀔 때마다 numpy의 너그러운 규칙으로 다시 계산한다.
\end{instardust}

\section{원소별 연산의 반복}
\label{sec:nd-iteration}

$d$차원 배열들의 원소별 연산을 가장 단순하게 쓰면 $d$겹의 반복문이다. 그러나 13장에서 보았듯이 메모리 배치와 어긋난 순서로 건너뛰며 훑으면 몇 배 느려지고, 가장 안쪽 반복이 짧으면 반복의 부담이 계산보다 커진다. 그래서 원소별 연산 엔진은 $d$차원 문제를 되도록 적고, 길고, 메모리를 따라가는 $1$차원 구간들로 바꾼 뒤, 그 구간마다 $1$차원 커널을 부른다. numpy의 \code{nditer}가 하는 일이고, 단계는 다음과 같다.

첫째로 모든 피연산자를 공통 모양으로 브로드캐스트하고, 둘째로 크기가 $1$인 축을 없앤다. 셋째로 축을 다시 배열해 보폭의 절댓값이 큰 축을 바깥에, 작은 축을 안쪽에 둔다. 출력의 보폭을 먼저 보아 쓰기가 차례로 이루어지게 하고, 보폭이 $0$인 축은 순서에 대한 정보가 없으므로 비교에서 뺀다. 넷째로 이웃한 두 축을 하나로 합칠 수 있으면 합친다. 다음 명제가 그 조건이다.

\begin{proposition}
\label{prop:coalesce}
바깥 축의 크기와 보폭이 $(n_o,s_o)$, 바로 안쪽 축이 $(n_i,s_i)$이고 모든 피연산자에서 $s_o=n_is_i$이면, 두 축을 크기 $n_on_i$, 보폭 $s_i$인 축 하나로 바꾸어도 모든 피연산자에서 원소의 주소와 방문 순서가 같다.
\end{proposition}

\begin{proof}
두 축의 첨자 $(j,k)$, $0\le j<n_o$, $0\le k<n_i$는 원래 반복에서 $j$가 바깥이므로 $t=jn_i+k$의 순서로 방문된다. 그 원소의 주소 기여는 $js_o+ks_i=jn_is_i+ks_i=(jn_i+k)s_i=ts_i$이다. 따라서 합친 축의 첨자 $t$가 $0$부터 $n_on_i-1$까지 차례로 갈 때 같은 주소를 같은 순서로 방문한다. 이 등식은 피연산자마다 따로 성립해야 하므로 조건은 모든 피연산자에 대해 필요하다.
\end{proof}

연속인 $(3,4,5)$ 배열들은 이 과정을 거쳐 길이 $60$인 구간 하나가 되고, 전치된 배열은 다시 배열하는 단계에서 메모리 순서로 읽힌다. 축들이 서로 다른 순서로 놓인 피연산자들을 섞으면, 예컨대 C 순서 배열에 그 전치를 더하면, 어느 축 순서를 골라도 한 피연산자는 메모리를 건너뛰며 읽게 된다. 이때는 13장의 블록화처럼 작은 타일 단위로 돌아 건너뛰며 읽는 쪽의 캐시 줄을 다시 쓰게 한다.

\begin{instardust}[\code{docs/design/array.md} §4.2]
stardust의 \code{NdIter}가 위의 네 단계를 하고, 바깥 축들은 자동차의 주행 거리계처럼 첨자를 올려 가며 돌면서 가장 안쪽 구간마다 커널을 한 번 부른다. 피연산자의 안쪽 보폭이 바깥 보폭보다 길어서 메모리 순서가 섞인 경우에는 $16$행 $\times$ $128$열의 타일로 도는데, 타일의 행들이 그 피연산자의 캐시 줄($128$바이트에 원소 $16$개)을 다시 쓰고 $128$줄($16$KB)이 L1에 남는다. 설계 문서는 $1000\times1000$인 \code{x.T + y}가 $552\mu$s에서 $430\mu$s로 줄었다고 적는다. 타일 순회는 원소마다 따로 계산하는 연산에만 쓰고 리덕션에는 쓰지 않는다. 방문 순서가 바뀌면 합하는 순서가 바뀌어 비트가 달라지기 때문이다(14장).
\end{instardust}

\section{자료형과 승격}
\label{sec:promotion}

서로 다른 자료형의 피연산자를 섞으면 어느 자료형으로 계산할지를 정해야 한다. 자료형들 사이에 ``값을 잃지 않고 바꿀 수 있다''는 부분순서를 두고 두 피연산자의 상한을 결과의 자료형으로 삼는 것이 표준적인 방법이며, 이 상한을 구하는 규칙을 \term{승격}{promotion}이라 한다. 자료형이 \code{bool}, \code{int64}, \code{float64}, \code{complex128} 넷뿐이면 이 순서는 전순서 $\code{bool}<\code{int64}<\code{float64}<\code{complex128}$로 둘 수 있고, 승격은 두 자료형의 최댓값이다. 엄밀히는 \code{int64}에서 \code{float64}로 가는 변환이 8장에서 보았듯이 $2^{53}$보다 큰 정수의 끝자리를 잃지만, 크기의 범위를 잃지 않는다는 뜻에서 위로 간다고 본다.

연산마다 모든 자료형의 반복문이 있는 것은 아니다. 산술 연산에는 \code{bool} 반복문이 없어 \code{True + True}는 \code{int64}로 계산되어 $2$이고, 나눗셈과 초월함수에는 정수 반복문이 없어 정수는 \code{float64}로 계산된다. 그래서 실제 규칙은 ``승격한 자료형부터 위로 올라가며 그 연산의 반복문이 있는 첫 자료형''이 된다. 자료형이 다른 입력은 계산 전에 변환해야 하는데, 입력 전체를 미리 변환한 임시 배열을 만들면 메모리를 두 번 더 지나므로, 작은 조각씩 변환해 바로 계산하는 \term{버퍼 단위 변환}{buffered casting}이 더 빠르다. 조각의 크기는 피연산자들의 조각이 함께 L1 캐시에 들어가도록 정한다.

\begin{instardust}[\code{docs/design/array.md} §2.1, §4.3]
stardust의 배열 자료형은 위의 넷이고 승격은 \code{promote(a, b) = max(a, b)} 한 줄이다. Python의 \code{int}, \code{float}, \code{complex} 스칼라도 같은 규칙을 따르며, numpy의 ``약한 스칼라'' 같은 예외는 두지 않는다. 버퍼 단위 변환의 조각은 원소 $4096$개인데, 피연산자 셋의 조각이 각각 $32$KB여서 성능 코어의 L1 캐시 $128$KB에 함께 들어간다. 설계 문서는 $256$부터 $65536$까지 재어 $1024$에서 $4096$이 가장 빨랐고 $16384$ 이상은 L1을 넘쳐 $10$--$25\%$ 느렸다고 기록한다.
\end{instardust}

\section{메모리 겹침}
\label{sec:overlap}

뷰가 있으면 입력과 출력이 같은 메모리를 공유할 수 있다. \code{a[1:] += a[:-1]}을 생각하자. 앞에서부터 차례로 계산하면, $a_2$를 갱신할 때 읽는 $a_1$은 이미 갱신된 값이다. numpy와 같은 라이브러리는 이런 경우에도 ``모든 입력을 다 읽은 뒤 출력을 쓴 것''과 같은 결과를 약속하므로, 입력과 출력이 겹치는지를 판정해 겹치면 입력을 먼저 복사해야 한다.

정확한 판정은 어렵다. 두 뷰가 같은 바이트를 공유하는지는 보폭들의 정수 일차결합이 같아지는지를 묻는 정수 방정식의 문제여서, 일반적으로는 빠르게 풀 수 없다. 그래서 보수적인 판정을 쓴다. 먼저 두 뷰가 닿을 수 있는 바이트 범위 $[\min,\max]$가 겹치지 않으면 공유하지 않는다. 겹치면 2장에서 본 최대공약수 판정을 쓴다. 모든 보폭의 최대공약수를 $g$라 하면 두 뷰의 원소 주소는 각각 시작 주소에 $g$의 배수를 더한 것이므로, 두 시작 주소를 $g$로 나눈 나머지가 원소 크기만큼의 틈을 두고 엇갈려 있으면 어떤 바이트도 겹칠 수 없다. \code{x[::2]}와 \code{x[1::2]}가 그 예이다. 이 판정은 충분조건이므로 확실히 겹치지 않는 경우만 가려내고, 나머지는 겹친다고 보고 복사한다.

\section{열 지향 표}
\label{sec:columnar}

\subsection{열마다 따로}

표 형식의 데이터는 행 단위로 저장할 수도 있고 열 단위로 저장할 수도 있다. 데이터 분석은 대개 한 열의 모든 값을 훑는 연산, 예컨대 합이나 평균이나 필터이므로, 한 열의 값들이 메모리에 연속으로 놓인 \term{열 지향}{columnar} 형식이 캐시와 SIMD에 유리하다. 열마다 자료형이 다르므로 표 전체를 하나의 $2$차원 배열로 두지 않고 열마다 따로 $1$차원 버퍼를 둔다. Apache Arrow가 이 형식을 언어와 라이브러리 사이에서 공유하는 표준으로 정했고, 오늘날의 데이터프레임 라이브러리 대부분이 그 형식을 따르거나 참고한다\cite{arrow-columnar}.

\subsection{결측과 비트맵}

데이터 분석의 표에는 값이 없는 칸, 곧 \term{결측}{missing value}이 있다. 부동소수점 열이라면 NaN으로 나타낼 수도 있지만, 정수와 문자열과 불 열에는 그런 특수값이 없고, NaN은 계산의 결과로도 생기므로 ``값이 없음''과 구별해야 할 때가 있다. 열 지향 형식은 값의 버퍼와 따로 행마다 $1$비트씩의 \term{유효 비트맵}{validity bitmap}을 두어, 비트가 $1$이면 값이 있고 $0$이면 결측으로 나타낸다. Arrow의 규칙에서 행 $i$는 바이트 $\lfloor i/8\rfloor$의 비트 $i\bmod8$이고(낮은 비트가 먼저), 결측이 하나도 없으면 비트맵 자체를 두지 않는다. 그러면 결측이 없는 흔한 경우에 커널은 비트맵을 아예 보지 않는 빠른 반복을 고를 수 있다.

결측인 칸의 값 버퍼에는 무엇을 둘까. 아무 값이나 두면 해시와 비교에 쓰레기 값이 섞여 결과가 실행마다 달라질 수 있으므로, 언제나 $0$을 두는 것이 좋다. 그러면 결측이 있는 부동소수점 열의 합을 분기 없이 모든 칸을 더해서 구할 수 있다. 결측 칸이 $0.0$을 더하기 때문이다.

\subsection{문자열}

길이가 제각각인 문자열은 한 칸에 담을 수 없으므로, 모든 문자열의 UTF-8 바이트(12장)를 이어 붙인 버퍼 하나와, 행 $i$의 문자열이 그 버퍼의 어디에서 시작하고 끝나는지를 담은 \term{위치표}{offsets} $o_0\le o_1\le\cdots\le o_n$을 둔다. 행 $i$의 문자열은 바이트 구간 $[o_i,o_{i+1})$이다. 이 형식에서 해시와 비교와 필터는 Python 문자열 객체를 거치지 않고 바이트를 바로 읽는다. 대신 Python으로 값을 꺼낼 때마다 문자열 객체를 새로 만들어야 한다.

\begin{instardust}[\code{docs/design/frame.md} §4.3--§4.6]
stardust의 \code{Column}은 자료형, 길이, 시작 위치, 결측의 개수와 참조 계수 버퍼 셋, 곧 값, 유효 비트맵, 문자열의 위치표로 이루어진다. 잘라 내기는 시작 위치와 길이만 바꾸고 버퍼를 공유하며, 그래서 문자열 위치표는 잘라 낸 열에서도 다시 계산하지 않는 절대 위치이다. 위치표는 $64$비트 정수여서 한 열이 $2$GB를 넘을 수 있다(Arrow의 \code{large\_string}과 같다). 결측의 개수를 함께 들고 다니므로 ``결측이 없으면 비트맵이 없다''가 잘라 낸 열에서도 성립하고, 결측 칸의 값은 언제나 $0$이다. 비트맵을 행마다 비트로 꺼내 분기하면 느리므로, 비트맵 $1$바이트를 $256$칸 표로 읽어 행마다 $0$ 또는 $1$인 바이트 $8$개로 한꺼번에 펼친 뒤 분기 없이 고른다. 설계 문서에 따르면 이것으로 \code{isna}가 약 $8$배 빨라졌다.
\end{instardust}

\section{Copy-on-Write}
\label{sec:cow}

여러 객체가 같은 데이터를 공유하면 복사를 아낄 수 있지만, 한 객체를 고칠 때 다른 객체까지 바뀌는 문제가 생긴다. 데이터프레임에서 열 하나를 꺼내 \code{s = df["a"]}로 이름을 붙인 뒤 \code{s}의 값 하나를 바꾸면, 사용자는 \code{df}가 그대로이기를 기대할 것이다. 꺼낼 때마다 복사하면 안전하지만 느리다.

\term{Copy-on-Write}{copy-on-write}는 복사를 쓰기가 일어날 때까지 미룬다. 꺼내기와 잘라 내기는 버퍼를 공유하는 새 객체만 만들고, 어떤 객체에 쓰기가 일어나면 그때 버퍼를 함께 쓰는 다른 주인이 있는지 확인해서 있으면 먼저 복사한 뒤 쓴다. 주인의 수는 버퍼의 참조 계수로 알 수 있다. 참조 계수가 $1$이면 아무도 함께 쓰지 않으므로 그 자리에서 고쳐도 된다. 운영체제가 프로세스를 복제할 때 메모리 페이지를 나누어 주는 방식과 같은 생각이다.

이 규칙이 올바르려면 버퍼를 가리키는 모든 길이 참조 계수에 잡혀 있어야 한다. 그래서 외부의 배열과 메모리를 공유하는 것은 위험하다. numpy 배열은 다른 쓰기 가능한 뷰를 통해 언제든 바뀔 수 있는데, 그 쓰기는 참조 계수에 드러나지 않기 때문이다. 그래서 바깥에서 들어오는 데이터는 복사하고, 바깥으로 내보내는 뷰는 읽기 전용으로 만든다.

\begin{instardust}[\code{docs/design/frame.md} §5.2, §5.4]
stardust의 데이터프레임에서 쓰기는 두 단계이다. 쓰는 쪽은 먼저 자기만의 \code{Column} 객체를 갖고, 행에 쓰는 함수가 값과 비트맵 버퍼의 참조 계수를 본다. $1$이면 그 자리에서 쓰고, 다른 주인이 있으면 자기 행 범위만 새 버퍼로 복사한 뒤 쓴다. 문자열 열은 값의 길이가 바뀔 수 있으므로 늘 새 버퍼를 만든다. 열을 배열로 내보낼 때는 같은 버퍼를 가리키는 읽기 전용 뷰를 주는데, 그 뷰가 살아 있는 동안 참조 계수가 하나 늘어 있으므로 데이터프레임 쪽의 쓰기는 위의 규칙대로 먼저 복사한다. 배열 쪽에서 들어오는 데이터는 언제나 복사한다. 여러 열을 하나의 $(n,k)$ 배열로 내보낼 때는 F 순서로 만드는데, 그러면 열 하나가 연속된 한 덩어리여서 열마다 \code{memcpy} 한 번이면 되고, 선형대수 영역이 가장 빠르게 읽는 배치도 F 순서이다.
\end{instardust}

\section{맺음}

배열과 표의 메모리 모델은 같은 생각의 두 변형이다. 데이터 자체와 그것을 읽는 규칙을 나누어 두면, 규칙만 바꾸는 연산은 공짜가 되고 데이터를 옮기는 연산만 비용을 낸다. 배열에서 그 규칙은 모양과 보폭이고, 표에서는 열의 시작 위치와 길이와 비트맵이다. 다음 장에서는 이 데이터 구조를 Python에서 C++로, 다시 C++에서 Python으로 넘기는 언어의 경계를 다룬다.

\chapterreferences
